\documentclass[11pt]{article}
\usepackage[margin=1in]{geometry}
\usepackage{amsmath}
\usepackage{xcolor}
\usepackage{parskip}
\usepackage{hyperref}
\usepackage{microtype}

\newenvironment{reviewer}{\color{blue!50!black}\itshape}{\normalcolor}
\newcommand{\response}{\textbf{Response:} }

\begin{document}

\begin{center}
{\Large \textbf{Response to Reviewers --- Second Round}}\\[6pt]
Manuscript jrfm-4438773\\
\emph{Selective Prediction and the Persistence Illusion: A Diagnostic Decomposition of VIX
Regime Classification}\\[4pt]
Akshat Gupta and Jianguo Liu
\end{center}

We thank both reviewers for their positive assessments and for the further suggestions, all of
which we have implemented. No numerical results changed in this round; the edits are editorial
additions and clarifications as detailed below.

\section*{Response to Reviewer 1}

We are grateful for the reviewer's generous summary of the first revision and for the three
remaining suggestions.

\textbf{Comment 1 (Results sections could be shorter; numerical discussion duplicated by
tables).}

\response We have trimmed the Results prose that restated table values. In particular, the
main-results discussion (Section 5.1) and the persistence-quantification discussion
(Section 5.4) now interpret the tables rather than repeat their entries, referring the reader
to the tables for the numbers. We note that the larger structural version of this
change was made in the first revision, when the threshold-sensitivity, VIX-threshold
robustness, and sustained-label tables were moved to Appendix B; the present pass completes it
at the paragraph level.

\textbf{Comment 2 (captions should remind readers which accuracy is reported).}

\response We audited every figure and table caption and added a one-line clarifier wherever
the base population was not already stated. Seven captions were amended (Figures 1, 3, and 6
and Tables 7 and 18--20) to state explicitly whether the reported quantity is
selective accuracy on covered days, baseline accuracy on all test days, balanced accuracy, or
an abstention rate over all days.

\textbf{Comment 3 (when is selective prediction still beneficial from a risk-management
perspective).}

\response We expanded the Practical Implications subsection with a dedicated risk-management
passage describing three concrete deployment patterns that survive the paper's negative
results: (i) \emph{abstention as an uncertainty signal}---treating a spike in the abstention
rate (which rises to 54\% immediately before calm$\to$high transitions at $N = 5$) as a
trigger for hedge review, position-limit tightening, or de-risking, so that the system's
silence rather than its directional calls carries the signal; (ii) \emph{regime-stable
confirmation}---using the 99--100\% covered accuracy on stable days to gate strategies whose
principal risk is false exit signals, such as systematic carry or volatility-selling programs;
and (iii) \emph{avoiding forced low-confidence calls}---routing abstained days to a human or a
conservative default in automated pipelines. The passage closes with the deployment principle
that selective systems should be used for what the diagnostic shows they can do (confirm
continuation, flag their own uncertainty) and not for what they cannot (anticipate spikes).

\section*{Response to Reviewer 2}

\textbf{Comment (add more citations to strengthen the paper).}

\response We have added five references, woven into the Introduction and Related Work where
they support the argument rather than appended as a list: Bekaert and Hoerova (2014,
\emph{Journal of Econometrics}) on the decomposition of the VIX into conditional variance and
the variance risk premium; Bloom (2009, \emph{Econometrica}) on uncertainty shocks as the
macroeconomic driver of volatility spikes; Andersen, Bollerslev, Diebold and Labys (2003,
\emph{Econometrica}) on realized volatility modeling; Paye (2012, \emph{Journal of Financial
Economics}) on the limited out-of-sample value of macroeconomic predictors for volatility,
which foreshadows our SHAP findings; and Gu, Kelly and Xiu (2020, \emph{Review of Financial
Studies}) as the benchmark for genuine machine-learning gains in empirical asset pricing,
against which we position our message as complementary: flexible models can add value, but on
highly persistent classification targets their apparent gains require persistence-matched
auditing.

\textbf{On the review-form item ``Is the research design appropriate?\ --- Must be
improved.''}

\response The reviewer's written comments raise no design concern, so we briefly address the
checkbox here for completeness. The research design was substantially strengthened in the
first revision in direct response to both reviewers and the editor: all baseline abstention
thresholds are now calibrated on the training window only and frozen before test evaluation
(eliminating the coverage-matching leakage identified in round one); the benchmark set was
extended with XGBoost and a continuous-forecast-then-threshold HAR model; evaluation was made
threshold-free via risk--coverage curves and AURC; inference robustness was documented across
bootstrap block lengths of 32--90 days; and the complete pipeline was published at
\url{https://github.com/akshatg10104/VIX_research} for independent verification. Reviewer 1,
who reviewed these changes in detail, assessed the design as appropriate in this round. We
would of course be glad to address any specific design concern the reviewer wishes to raise.

\vspace{1em}
We thank both reviewers again for reviews that have materially improved the paper.

\begin{flushleft}
Sincerely,\\[4pt]
Akshat Gupta and Jianguo Liu
\end{flushleft}

\end{document}
